\CharacterSheet{prevlover}{Antigo amante}{

\shortdescription{O investigador particular da família.}

\ObjectivesHeading{Objetivos}
\begin{itemize}
  \item Descobrir quem matou \Character{newcomer} e conseguir que essa pessoa seja expulsa por votação. (\PositiveReward{1}, \NegativeReward{1})
\end{itemize}

\LosingConditionsHeading{Condições de derrota}
Não conseguir expulsar por votação quem matou \Character{newcomer}, caso haja uma pessoa responsável pelo homicídio.

\SecretsHeading{Segredos}
\begin{itemize}
  \item Tive uma relação secreta com \Character{newcomer}. (\NegativeReward{1}, \PositiveReward{1})
  \item Usei o ritual roubado na memória de \Character{enforcer}. (\NegativeReward{1}, \PositiveReward{1})
\end{itemize}

\ActionsHeading{Ações}
\begin{itemize}
  \item Trata cada afirmação como uma pista. Pergunta quem viu, quando aconteceu e que provas restam.
\end{itemize}

\BackgroundHeading{Contexto}

\timestamp{No ano 1000 d.C.}, juntei-me à família porque era bom a reparar no que os outros preferiam esconder. Não sou de me gabar, mas resolvi crimes em pelo menos uma aldeia. \Character{vampirehead} queria um investigador e escolheu-me. Estava entusiasmado por me tornar um ser sobrenatural e adorei a festa que me fizeram. \Character{vampirehead}, \Character{usurper}, \Character{enforcer}, \Character{buttler}, \Character{ghoul}, \Character{supremacist} e \Character{lover} foram uma companhia surpreendentemente agradável.

\Character{usurper} deu-me um pequeno artefacto que ajudava a distinguir a verdade da mentira em afirmações escritas. Não me lembro bem dos pormenores do artefacto, apenas de que só funciona uma vez por dia.

\Character{enforcer} e eu trabalhávamos juntos com frequência. Embora não parecesse gostar particularmente de mim, achava-me útil.

\Character{traitor} chegou em \timestamp{1400 d.C.}, depois de \Character{lover} ter implorado que a admitissem. Tanto eu como os outros achávamos que a sua bondade destoava nesta casa.

\timestamp{Há dois anos}, \Character{vampirehead} apresentou \Character{newcomer} num encontro de família. O interesse da chefe por ele era evidente, tal como o perigo de me aproximar dele. Fi-lo na mesma. Começámos uma relação em segredo. Durante algum tempo, pensei ter encontrado alguém capaz de devolver novidade a uma vida imortal.

\timestamp{Há um ano}, \Character{newcomer} acabou a relação. Disse-me: ``O que estamos a fazer está errado'', sem acrescentar mais nada. Presumi que o interesse da chefe o tivesse assustado.

Por essa altura, \Character{enforcer} pediu-me ajuda para investigar uma página desaparecida do \Clue{rulebook}. Suspeitava de \Character{traitor}, em parte pela sua hostilidade à família. Intercetámos correspondência invulgar ligada às casas de \Character{traitor} e de \Character{lover}. As mensagens pareciam importantes, mas a investigação deixou mais pontas soltas do que conclusões. Guardei as notas, juntamente com o meu artefacto, na minha mala. O que estava lá escrito? Quem me dera lembrar-me. Juro que a minha memória costuma ser muito melhor.

\timestamp{Há dois meses}, \Character{vampirehead} anunciou que \Character{newcomer} se juntaria à família numa festa em \timestamp{16/7/2026}, e todos os vampiros tinham de comparecer.

\timestamp{Umas duas semanas antes da festa}, \Character{vampirehead} perguntou se mais alguém, além de \Character{usurper}, \Character{buttler} e \Character{newcomer}, queria chegar duas noites antes para ajudar. Recusei. O convite deixava claro que o gabinete de \Character{usurper} ficaria sem vigilância na manhã anterior à festa.

\timestamp{Ontem, à 01:00}, entrei no gabinete de \Character{usurper} à procura de uma forma de compreender \Character{newcomer} e talvez fazê-lo amar-me de novo. Encontrei um ritual chamado ``Mudar e conhecer uma pessoa''. Prometia acesso às memórias de uma pessoa e o poder de as alterar. Roubei-o. Alguém me viu no gabinete. Na pressa de sair, deixei cair a mala, com o artefacto que permite ver a verdade e as notas sobre a correspondência. Ainda não recuperei nada.

\timestamp{Ontem, às 16:00}, \Character{enforcer} telefonou-me. Disse que ignoraria o roubo se eu lhe entregasse o ritual. Queria usá-lo em \Character{newcomer} porque alguém a tinha avisado de que ele morreria. Percebi que examinar as memórias dele exporia a nossa relação, e a família teria motivos para me retirar os poderes. Entrei em pânico.

\timestamp{Às 16:15}, usei o ritual na memória de \Character{enforcer} antes que ela o pudesse usar na de \Character{newcomer}. Queria fazê-la esquecer o acordo, mas apressei o ritual e algo correu mal. Sei-o porque foi nesse preciso momento que também comecei a ter problemas de memória. Parece que nos afetei aos dois. O ritual produziu resíduos.

\timestamp{Às 17:15}, reparei que os resíduos tinham desaparecido. Onde estarão?

\GMNotes{
\paragraph{Porque se pode suspeitar desta personagem?}
\Character{prevlover} tinha um motivo amoroso doloroso, entrou no gabinete de \Character{usurper} sem autorização, roubou magia poderosa sobre a memória e apagou deliberadamente a memória de \Character{enforcer}. O artefacto e as notas sobre as mensagens intercetadas, ambos desaparecidos, também podem dar a impressão de que destruiu ou escondeu provas.

\paragraph{Como se pode provar a inocência desta personagem?}
O crime de \Character{prevlover} ocorreu às \timestamp{16:15}, mais de quatro horas depois de \Character{newcomer} ter bebido o vinho de \Character{usurper} ao almoço. O ritual da memória causou confusão e ocultou provas, mas não colocou a primeira dose no vinho. O telefonema de \Character{enforcer} e o aviso anterior de \Character{ghoul} demonstram que \Character{prevlover} reagia a uma morte que já estava em marcha.

\paragraph{Qual é o papel desta personagem na história e no jogo?}
\Character{prevlover} é simultaneamente um investigador comprometido e uma via para várias pistas físicas. A mala que deixou cair explica por que razão o artefacto que permite ver a verdade e os registos da correspondência podem ser encontrados longe do dono. O que sabe sobre a tentativa de intervenção de \Character{enforcer} faz do ritual da memória um crime relevante, sem o tornar no homicídio.

\paragraph{Como pode esta personagem encaminhar a investigação para a solução?}
Assim que admitir o que \Character{enforcer} disse ao telefone, o grupo terá motivos para apurar quando foi prevista a morte pela primeira vez. \Character{ghoul} poderá então situar o aviso às \timestamp{15:00}, depois do almoço em que \Character{usurper} participou. Essa ordem exclui os rituais posteriores como causa e leva a investigação de volta à troca do vinho ao almoço.
}
}
